\documentclass[10pt,a4paper]{amsart}
\usepackage[margin=19mm]{geometry}
\usepackage{amsmath,amssymb,mathtools}
\usepackage[scr=rsfs,cal=euler]{mathalfa}
\usepackage{xfrac}
\usepackage[unicode,hidelinks,bookmarks=false]{hyperref}
\newcommand{\ie}{i.e.\ }
\newcommand{\cf}{cf.\ }
\newcommand{\resp}{respectively\ }
\newcommand{\Subsection}{\subsection}
\providecommand{\nobreakdash}{\nobreak\hbox{-}}
\setlength{\emergencystretch}{2em}
\multlinegap0pt
\usepackage{amssymb}
\usepackage{xcolor}
\newtheorem{lem}{Lemma}
\newcommand{\Z}{\mathbb{Z}}
\newcommand{\cF}{\mathcal{F}}
\title{The Poisson Matrix $\mathbf{A}_2$ characteristic and the 3/2 blow up of the Hilbert transform}
\author{Komla Domelevo, Spyridon Kakaroumpas, Stefanie Petermichl, Sergei Treil, Alexander Volberg}
\date{Source: arXiv:2605.19637v1 (CC BY 4.0)}
\begin{document}
\maketitle

\noindent\emph{Source and licence.} The Poisson Matrix $\mathbf{A}_2$ characteristic and the 3/2 blow up of the Hilbert transform. Version arXiv:2605.19637v1 (CC BY 4.0), distributed by arXiv under the Creative Commons Attribution 4.0 International licence, http://creativecommons.org/licenses/by/4.0/, as recorded in the arXiv licence metadata for this exact version and shown on the abstract page https://arxiv.org/abs/2605.19637v1. Full licence notice: \texttt{licenses/arxiv-2605.19637-CC-BY-4.0.txt}.\par\smallskip

\noindent\textit{Editorial scope.} Counted unit: the article's lem lem:finiteness\_hitting\_times (source0.tex lines 501--521). The article's complete original proof of it is delivered with it (source0.tex lines 523--576, one \texttt{proof} environment of 54 lines). No mathematics is written, repaired or re-derived: the statement and the proof are the article's own lines, in the article's own notation, and no retained line is substituted or re-flowed.  No further source-local context is needed: every cross-reference of every retained line resolves inside the two delivered spans.  Nothing was omitted from any retained span, so the package carries no omission record.  No retained line contains a citation, so the wrapper carries no bibliography and no bibliography entry.  No retained line contains a figure environment, an image-inclusion command or drawing code, so no image is delivered and the package declares no asset dependency.  Editorial formatting only, in addition: an AMS \texttt{amsart} preamble that restates the article's own declarations and macros used by the retained lines, namely the environments \texttt{lem} on separate counters, together with the standard packages those lines need.  The retained environments print in this document as Lemma 1 in order of appearance, whereas the article prints the same items with the numbering of its own full document.  The complete native source of the article as distributed in the arXiv e-print for this exact version is bundled unchanged as \texttt{sources/200932/source0.tex}.
\begin{lem}
\label{lem:finiteness_hitting_times}
Let $(\Omega,\cF,\mathbb{P})$ be a probability space. Let $\{\omega_n\}^{\infty}_{n=1}$ be a sequence of $\Z$-valued, independent, identically distributed, not $\mathbb{P}$-a.s.~constant random variables on $\Omega$ with $\mathbb{E}[\omega_1]=0$. Let $a\in\Z$. Consider the sequence $\{S_n\}_{n=0}^{\infty}$ of random variables on $\Omega$ given by
\begin{equation*}
	S_0\equiv a,\quad S_n:=a+\sum_{k=1}^{n}\omega_k,~n=1,2,\ldots.
\end{equation*}

\begin{enumerate}
    \item Assume that there exists $M\in\Z$ with $\omega_n\leq M$ $\mathbb{P}$-a.s., for every $n=1,2,\ldots$. Let $b\in\Z$ with $a<b$. Set
    \begin{equation*}
	    \tau:=\inf\{n\geq1:~S_n\geq b\}.
    \end{equation*}
    Then, we have $\tau<\infty$ $\mathbb{P}$-a.s. Moreover, if $\omega_n\leq 1$ $\mathbb{P}$-a.s., for every $n=1,2,\ldots$, then $S_{\tau}=b$ $\mathbb{P}$-a.s.

    \item Assume that there exists $m\in\Z$ with $\omega_n\geq m$ $\mathbb{P}$-a.s., for every $n=1,2,\ldots$. Let $c\in\Z$ with $a>c$. Set
    \begin{equation*}
	    \sigma:=\inf\{n\geq1:~S_n\leq c\}.
    \end{equation*}
    Then, we have $\sigma<\infty$ $\mathbb{P}$-a.s. Moreover, if $\omega_n\geq 1$ $\mathbb{P}$-a.s., for every $n=1,2,\ldots$, then $S_{\sigma}=c$ $\mathbb{P}$-a.s.
\end{enumerate}
\end{lem}

\begin{proof}
\begin{enumerate}
    \item First of all, let $\cF_0:=\{\emptyset,\Omega\}$, $\cF_n:=\sigma(\omega_1,\ldots,\omega_n)$, $n=1,2,\ldots$ and consider the filtration $\mathbb{F}:=\{\mathcal{F}_n\}^{\infty}_{n=1}$ on $\Omega$. Fix a positive real number $\theta$ and set
	\begin{equation*}
		f(\theta):=\mathbb{E}[e^{\theta\omega_1}].
	\end{equation*}
	Then, it is easy to verify that the stochastic process $\{M_n\}^{\infty}_{n=0}$ given by
	\begin{equation*}
		M_n:=\frac{e^{\theta S_n}}{(f(\theta))^n}\cdot e^{-\theta a},\quad n=0,1,2\ldots
	\end{equation*}
	is a $\mathbb{F}$-adapted martingale. Obviously, $\tau$ is a $\mathbb{F}$-stopping time. Therefore, from the optional stopping theorem we deduce that the stopped process $\{M_{\tau\wedge n}\}^{\infty}_{n=0}$ is also a $\mathbb{F}$-adapted martingale. Let us note that for $\mathbb{P}$-almost every $x\in\Omega$ with $\tau(x)<\infty$, we have $\tau(x)>0$ and $S_{\tau(x)-1}(x)\leq b$, therefore $S_{\tau(x)}(x)\leq b+M$. Thus, we deduce
	\begin{equation*}
		0<M_{\tau\wedge n}\leq\frac{e^{\theta\max\{b,b+M\}}}{(f(\theta))^n}\cdot e^{-\theta a}\quad\mathbb{P}\text{-a.s.,}
	\end{equation*}
    for all $n=0,1,2,\ldots$. In particular, $\{M_{\tau\wedge n}\}^{\infty}_{n=0}$ is uniformly integrable. Therefore, there is a random variable $M$, such that $M_{\tau\wedge n}\rightarrow M$ pointwise $\mathbb{P}$-a.s.~and in $L^1(\Omega)$ as $n\rightarrow\infty$. 
	
	Notice that by Jensen's inequality we have
	\begin{equation*}
		f(\theta)>e^{\mathbb{E}[\theta \omega_1]}=1,
	\end{equation*}
	the strict inequality being due to the fact that $\omega_1$ is not $\mathbb{P}$-a.s.~constant. Thus, for all $x\in\Omega$ with $\tau(x)=\infty$, since
	\begin{equation*}
		0\leq M_{\tau(x)\wedge n}(x)=\frac{e^{\theta S_{n}(x)}}{(f(\theta))^n}\cdot e^{-\theta a}\leq\frac{e^{\theta b}}{(f(\theta))^n}\cdot e^{-\theta a},\quad\forall n=0,1,2,\ldots,
	\end{equation*}
	we deduce
	\begin{equation*}
		\lim_{n\rightarrow\infty}M_{\tau(x)\wedge n}(x)=0.
	\end{equation*}
	It follows that $M(x)=0$ for $\mathbb{P}$-almost every $x\in\Omega$ with $\tau(x)=\infty$. Observe also that
	\begin{equation*}
		M(x)=\frac{e^{\theta S_{\tau(x)}}}{(f(\theta))^{\tau(x)}}\cdot e^{-\theta a}
	\end{equation*}
	(because $\tau(x)\wedge n=\tau(x)$, for all $n=\tau(x),\tau(x)+1,\ldots$), for $\mathbb{P}$-almost every $x\in\Omega$ with $\tau(x)<\infty$. Therefore, we obtain
	\begin{equation*}
		\mathbb{E}[M]=\mathbb{E}\left[\frac{e^{\theta S_{\tau}}}{(f(\theta))^{\tau}}\cdot e^{-\theta a}\mathbf{1}_{\{\tau<\infty\}}\right].
	\end{equation*}
	Since also $\mathbb{E}[M]=\mathbb{E}[M_0]=1$, we deduce
	\begin{equation*}
		\mathbb{E}\left[\frac{e^{\theta (S_{\tau}-a)}}{(f(\theta))^{\tau}}\mathbf{1}_{\{\tau<\infty\}}\right]=1.
	\end{equation*}
	Since $\theta>0$ was arbitrary, letting $\theta\rightarrow0^{+}$ and using the Dominated Convergence Theorem we obtain $\mathbb{E}[\mathbf{1}_{\{\tau<\infty\}}]=1$, that is $\tau<\infty$ $\mathbb{P}$-a.s.

    If $\omega_n\leq 1$ $\mathbb{P}$-a.s., for every $n=1,2,\ldots$, then for $\mathbb{P}$-almost every $x\in\Omega$ with $\tau(x)<\infty$ we have $S_{\tau(x)-1}(x)\leq b-1$, therefore
	\begin{equation*}
		b\leq S_{\tau(x)}(x)\leq b-1+1=b,
	\end{equation*}
	thus $S_{\tau(x)}(x)=b$.

    \item This follows immediately from the first part after observing that $-a<-c$ and
    \begin{equation*}
		\sigma=\inf\{n\geq 1:~-S_n\geq-c\}.
	\end{equation*}
\end{enumerate}
\end{proof}

\end{document}
